%%%%%%%%%%%%%%%%%%%%%part-01.tex%%%%%%%%%%%%%%%%%%%%%%%%%%%%%%%%%%
% 
% sample part title
%
% Use this file as a template for your own input.
%
%%%%%%%%%%%%%%%%%%%%%%%% Springer %%%%%%%%%%%%%%%%%%%%%%%%%%

\begin{partbacktext}
\part{Morfología, geometría estructural y equivalencia topológica}

\noindent Esta segunda parte del documento establece la transición formal entre la teoría abstracta de los autómatas celulares, desarrollada en los Capítulos \ref{cap:preliminares} y \ref{cap:teoria_formal}, y su representación geométrica computable. El objetivo de esta sección es definir los modelos matemáticos y topológicos necesarios para proyectar la dinámica de una retícula hexagonal sobre la arquitectura ortogonal de la memoria de un equipo de cómputo convencional, garantizando en todo momento la rigurosidad analítica y la invariancia espacial.

En el Capítulo \ref{cap:geometria_computacional}, se formalizan las estructuras algebraicas que rigen el espacio discreto bidimensional. Se modela la retícula como un grupo abeliano bajo la adición y se definen las métricas analíticas fundamentales para la tesela hexagonal. De igual forma, se establece el modelo algebraico inverso para transformar el espacio continuo $\mathbb{R}^2$ a coordenadas discretas. Esto permite demostrar teóricamente que el algoritmo de actualización topológica opera en un tiempo asintóticamente lineal $\mathcal{O}(n)$ y con requerimientos espaciales de $\mathcal{O}(n)$, justificando empírica y formalmente el desarrollo del simulador en C/C++.

Posteriormente, en el Capítulo \ref{chap:isomorfismos}, se aborda el problema central de la representación en memoria: la transformación geométrica entre topologías. Se definen las ecuaciones de mapeo afín para proyectar la vecindad hexagonal pura hacia un tensor ortogonal de subred $2 \times 2$ (denotado como $\mathcal{T}_B$). A través de una demostración deductiva, se comprueba que la inactividad intencional de dos aristas en este bloque ortogonal previene la pérdida de simetría en la propagación de la información. 

En conjunto, el modelado algebraico y las demostraciones topológicas de esta parte constituyen el núcleo del motor matemático del proyecto. Este bloque garantiza estructuralmente que ninguna regla o propiedad de computabilidad se degrade al ser simulada, estableciendo el sustrato algorítmico indispensable para el análisis cinemático de partículas (\textit{gliders}) y el cálculo de interacciones locales booleanas que se detallarán en el Capítulo \ref{cap:dinamica_particulas}.

\end{partbacktext}